% Graphic paths are relative to resources folder
\graphicspath{{resources/}}

% TikZ - Definitions
\usepgfplotslibrary{fillbetween}
\pgfdeclarelayer{tree}
\pgfdeclarelayer{lines}
\pgfdeclarelayer{scope1}
\pgfdeclarelayer{scope2}
\pgfdeclarelayer{scope3}
\pgfplotsset{compat=1.15}
\usetikzlibrary{
	arrows.meta,
	automata,
	shapes,
	shapes.geometric,
	decorations,
	decorations.pathreplacing,
	angles,
	quotes,
	positioning,
	calc,
	shadows,
	trees,
	matrix,
	patterns,
	patterns.meta,
	calligraphy,
	fadings
}

% Special symbols
\newcommand{\checkyes}{\ding{51}~}

\newcommand{\correct}{\highlight{\ding{51}}}
\newcommand{\incorrect}{\highlight[secondary]{\ding{55}}}
\newcommand{\good}[1]{\textcolor{basebluegreen}{#1}}
\newcommand{\medium}[1]{\textcolor{orange}{#1}}
\newcommand{\bad}[1]{\textcolor{basered}{#1}}

% Married symbol, from https://tex.stackexchange.com/questions/180297/does-a-symbol-for-married-exist-in-some-package-or-font
\newcommand{\marriedto}[1][white]{\mathrel{\begin{tikzpicture}[x=1ex, y=1ex, scale=0.5, baseline=-.6ex]
			\begin{scope}
				\clip (-2,-.2) rectangle ++(5.5,2.2);
				\path[draw=#1,line width=.08ex,
					fill=black,even odd rule]
				(0, 0) circle (2ex)
				(0, 0) ellipse (1.5ex and 1.3ex);
				\path[draw=#1,line width=.08ex,
					fill=black,even odd rule]
				(1.5, 0) circle (2ex)
				(1.5, 0) ellipse (1.5ex and 1.3ex);
			\end{scope}
			\begin{scope}
				\clip (-2,-2) rectangle ++(5.5,2.2);
				\path[draw=#1,line width=.08ex,
					fill=black,even odd rule]
				(1.5, 0) circle (2ex)
				(1.5, 0) ellipse (1.5ex and 1.3ex);
				\path[draw=#1,line width=.08ex,
					fill=black,even odd rule]
				(0, 0) circle (2ex)
				(0, 0) ellipse (1.5ex and 1.3ex);
			\end{scope}
		\end{tikzpicture}}}

% Labeled trees in TikZ
\NewDocumentEnvironment{tikzltr}{O{1.5cm}O{}}{
\tikzpicture[sibling distance=#1, level distance=0.7cm, every node/.style={inner sep=3pt, minimum size=2mm}, edge from parent path={(\tikzparentnode) -- (\tikzchildnode)}, #2]
}{
\endtikzpicture
}


% QR code with link underneath
\newcommand*{\qrlink}[1]{{\qrcode{#1}\\[2ex]\url{#1}}}

% Source annotation for graphics
\newcommand{\source}[1]{%
	\caption*{\centering \fontsize{8}{0} \selectfont \thinspace{\url{#1}}}%
}
\newcommand{\tsource}[1]{%
	\caption*{\centering \fontsize{8}{0} \selectfont \thinspace{#1}}
}


% "Schleudern-Schild": Danger of vehicle skidding when road is wet or soiled (see German road traffic regulations [Straßenverkehrsordnung], sign no. 114)
\newcommand{\dangerOfVehicleSkiddingWhenRoadIsWetOrSoiled}{\includegraphics[scale=0.2]{common/slippery.pdf}}
\newcommand{\slippery}{\dangerOfVehicleSkiddingWhenRoadIsWetOrSoiled}


%******************************************************************************
% Magic which allows to cross out a frame
% taken from http://tex.stackexchange.com/a/109162
\usetikzlibrary{calc,decorations.pathmorphing,patterns}

\makeatletter

\pgfdeclaredecoration{penciline}{initial}{
	\state{initial}[width=+\pgfdecoratedinputsegmentremainingdistance,auto corner on length=1mm,]{
		\pgfpathcurveto%
		{% From
			\pgfqpoint{\pgfdecoratedinputsegmentremainingdistance}{\pgfdecorationsegmentamplitude}
		}
		{%  Control 1
			\pgfmathrand
			\pgfpointadd{\pgfqpoint{\pgfdecoratedinputsegmentremainingdistance}{10pt}}
			{\pgfqpoint{-\pgfdecorationsegmentaspect\pgfdecoratedinputsegmentremainingdistance}%
				{\pgfmathresult\pgfdecorationsegmentamplitude}
			}
		}
		{%TO
			\pgfpointadd{\pgfpointdecoratedinputsegmentlast}{\pgfpoint{8pt}{5pt}}
		}
	}
	\state{final}{}
}

%% Environment crossedOutFrame
% behaves like frame, but allows putting a big cross over the slide
% if option animate is given, the slide is crossed out on the last overlay
\newenvironment{crossedOutFrame}[2][]{
	\def\cofo{#1}
	\begin{frame}{#2}
		}{
		\IfSubStr{\cofo}{animate}{\pause}{ }
		\begin{tikzpicture}[remember picture,overlay,decoration=penciline]
			\draw[decorate,line width=30pt,red!60!black]
			(current page.north west) -- (current page.south east);
			\draw[decorate,line width=30pt,red!60!black]
			(current page.north east) -- (current page.center) -- (current page.south west);
		\end{tikzpicture}
	\end{frame}
}


% TikZ defaults
\tikzset{
	normal/.style={draw, semithick},
	n/.style={style=normal, circle, inner sep=1mm, minimum size=8mm},
	l/.style={style=normal, rounded corners=1mm, inner sep=1mm, minimum size=6mm},
	e/.style={style=normal, shorten >=1mm, shorten <=1mm, ->},
	syntax/.style={style=normal, ellipse, minimum height=6mm, minimum width=8mm}, % nodes in syntax trees
	inner/.style={style=normal, minimum size=4mm}, % inner leaves or root in normal trees
	leaf/.style={style=normal, circle, minimum size=4mm}, % leaves in normal trees
	te/.style={style=normal}, % edges in a tree
	be/.style={style=e, dashed} % binding edge
}

% TikZ syntax trees
\newenvironment{tikzsyntaxtree}[1][]{
	\begin{tikzpicture}[baseline=(current bounding box.north), #1]
		\tikzset{grow=down}
		\tikzset{every tree node/.style={syntax}}
		\tikzset{edge from parent/.style={te, edge from parent path={(\tikzparentnode) -- (\tikzchildnode)}}}
		}{
	\end{tikzpicture}
}

% Rectangular drop shadows
% code adapted from https://tex.stackexchange.com/a/11483/3954
\def\shadowshift{1pt,-1pt}
\def\shadowradius{8pt}
\colorlet{innercolor}{black!40}
\colorlet{outercolor}{gray!05}

% this draws a shadow under a rectangle node
\newcommand\drawshadow[1]{
	\begin{pgfonlayer}{shadow}
		\shade[outercolor,inner color=innercolor,outer color=outercolor] ($(#1.south west)+(\shadowshift)+(\shadowradius/2,\shadowradius/2)$) circle (\shadowradius);
		\shade[outercolor,inner color=innercolor,outer color=outercolor] ($(#1.north west)+(\shadowshift)+(\shadowradius/2,-\shadowradius/2)$) circle (\shadowradius);
		\shade[outercolor,inner color=innercolor,outer color=outercolor] ($(#1.south east)+(\shadowshift)+(-\shadowradius/2,\shadowradius/2)$) circle (\shadowradius);
		\shade[outercolor,inner color=innercolor,outer color=outercolor] ($(#1.north east)+(\shadowshift)+(-\shadowradius/2,-\shadowradius/2)$) circle (\shadowradius);
		\shade[top color=innercolor,bottom color=outercolor] ($(#1.south west)+(\shadowshift)+(\shadowradius/2,-\shadowradius/2)$) rectangle ($(#1.south east)+(\shadowshift)+(-\shadowradius/2,\shadowradius/2)$);
		\shade[left color=innercolor,right color=outercolor] ($(#1.south east)+(\shadowshift)+(-\shadowradius/2,\shadowradius/2)$) rectangle ($(#1.north east)+(\shadowshift)+(\shadowradius/2,-\shadowradius/2)$);
		\shade[bottom color=innercolor,top color=outercolor] ($(#1.north west)+(\shadowshift)+(\shadowradius/2,-\shadowradius/2)$) rectangle ($(#1.north east)+(\shadowshift)+(-\shadowradius/2,\shadowradius/2)$);
		\shade[outercolor,right color=innercolor,left color=outercolor] ($(#1.south west)+(\shadowshift)+(-\shadowradius/2,\shadowradius/2)$) rectangle ($(#1.north west)+(\shadowshift)+(\shadowradius/2,-\shadowradius/2)$);
		\filldraw ($(#1.south west)+(\shadowshift)+(\shadowradius/2,\shadowradius/2)$) rectangle ($(#1.north east)+(\shadowshift)-(\shadowradius/2,\shadowradius/2)$);
	\end{pgfonlayer}
}

% create a shadow layer, so that we don't need to worry about overdrawing other things
\pgfdeclarelayer{shadow}
\pgfsetlayers{shadow,main}
\newsavebox\shadowbox
\newlength\shadowlen
\newcommand\shadowimage[2][]{%
	\setbox\shadowbox=\hbox{\includegraphics[#1]{#2}}
	\setlength\shadowlen{\wd\shadowbox}
	\ifnum\shadowlen<\ht\shadowbox
		\setlength\shadowlen{\ht\shadowbox}
	\fi
	\divide \shadowlen by 120
	\def\shadowshift{\shadowlen,-\shadowlen}
	\def\shadowradius{\the\dimexpr\shadowlen+\shadowlen+\shadowlen\relax}
	\begin{tikzpicture}
		\node[anchor=south west,inner sep=0] (image) at (0,0) {\includegraphics[#1]{#2}};
		\drawshadow{image}
	\end{tikzpicture}%
}
